\section{Related Work}
\label{sec:related-work}

\paragraph{Railway disruption prediction.} Operational railway analytics has historically focused on short-horizon delay propagation models calibrated on a single operator's data~\cite{gerum2024railway}. Single-country studies dominate the literature and the data harmonisation problem is rarely addressed: the same nominal feature (e.g.\ \emph{weather severity}) carries different semantics across operators because the underlying observations come from different agencies (Trenitalia internal weather labels, FMI for Finland, Open-Meteo Historical Archive for the Netherlands~\cite{zippenfenig2023openmeteo}). Our pipeline puts heavy engineering effort into an anti-leakage Common Data Model that resolves these mismatches before any model is trained.

\paragraph{Tabular gradient-boosted ensembles.} For spatio-temporal tabular problems the established baselines are LightGBM~\cite{ke2017lightgbm} and XGBoost~\cite{chen2016xgboost} -- both gradient-boosted decision-tree systems with histogram-based split-finding that scale linearly to tens of millions of rows. Their dominance over deep models on heterogeneous tabular data is well documented; we treat them as default baselines and we add a Random Forest~\cite{breiman2001random} for variance comparison.

\paragraph{Graph and sequence models.} We also benchmark a graph neural network -- GraphSAGE~\cite{hamilton2017inductive} -- to test whether the station-adjacency graph (163,902 directed edges across 2,397 stations) carries signal beyond hand-engineered station-static features (degree, betweenness, average historical delay). The bidirectional / causal LSTM~\cite{schuster1997bidirectional,hochreiter1997long} variant is an obvious next step but was deferred for memory reasons (see \cref{sec:models}).

\paragraph{Explainability.} We rely on TreeSHAP~\cite{lundberg2017unified,lundberg2020local}, the canonical exact attribution method for tree ensembles, and run it on the \emph{full} 2.30~M-row test set rather than a sample. SHAP gives us the global feature importance and the per-row local attributions that the Knowledge Graph bridge query (\cref{sec:kg}) consumes.

\paragraph{Graph databases.} Sparksee~\cite{sparksee2024} (descendant of DEX~\cite{martinez2017dex}) was the original target backend for the Knowledge Graph -- the project ships a valid Sparksee licence file. Because the binary distribution is gated behind a registered-user login that the build environment cannot reach, we instead use Kuzu~\cite{feng2023kuzu}, a modern Cypher-native embedded graph database that supports identical schema definitions and bulk-load Cypher commands. The schema, the staged parquets and the bridge query are all 100\,\% portable to Sparksee whenever its binary is available.

\paragraph{Cross-country transfer in operational ML.} Transfer learning across operators is a recurring practical challenge: the source-domain feature distributions rarely match the target. We document one such failure mode in \cref{sec:transfer} where a Netherlands-trained cause model fails on Finland because of incompatible precipitation and temperature scales between Open-Meteo daily aggregates and FMI per-stop hourly observations. Our negative-result documentation contributes to the broader practical literature on when domain-naive transfer is unsafe.
